\documentclass[11pt,a4paper]{article}
\usepackage[T1]{fontenc}
\usepackage{amsmath,amssymb,geometry,booktabs,parskip,xcolor,hyperref}
\geometry{margin=2.3cm}
\hypersetup{colorlinks=true,linkcolor=blue!60!black,urlcolor=blue!60!black}
\title{Stage 2: derive the execution HJB and classical solution}
\author{Equations, derivations and assumptions}
\date{}
\begin{document}
\maketitle
% Self-contained for Overleaf: this is a copy of derivation.tex, which
% WorkPlan.tex also \input's. Re-merge the two if derivation.tex changes.
\setcounter{section}{1}
\section{Step 2: specialize the 2022 paper's Eq.~16}
\label{sec:stage2}

\subsection*{Main equation and what this stage establishes}

The unknown $\mathcal V(\mathbf x,t)$ is the minimum remaining execution
cost for inventory $\mathbf x$ at time $t$. This stage derives
\[
 \boxed{\mathcal V_t-\tfrac14(\nabla\mathcal V)^\top H^{-1}\nabla\mathcal V
 +\lambda\mathbf x^\top C\mathbf x=0,\qquad
 \mathbf v^*=\tfrac12H^{-1}\nabla\mathcal V.}
\]
The first term measures how the cost changes with time, the second is
the optimized trading cost, and the last penalizes inventory risk.
We first reduce the source equation, then derive the same HJB from the
cost functional to check the signs and numerical factors independently.

\subsection*{Assumptions used in the derivation}

\begin{enumerate}
\item Inventory evolves continuously as $\dot{\mathbf x}=-\mathbf v$ on
a finite horizon $T>0$. Controls are square-integrable signed rates in
$\mathbb R^M$; purchases are allowed in this reference problem.
No nonnegativity constraint, spread or participation cap is imposed.
\item The impact matrix $H=\operatorname{diag}(\eta_j)$ is constant and
$\eta_j>0$. Thus $\mathbf v^\top H\mathbf v$ is strictly convex and
$H^{-1}$ exists. The constant symmetric covariance satisfies $C\succeq0$,
and $\lambda\ge0$.
\item In the source model, $R=0$ removes mean reversion and leaves
driftless Brownian prices. The terminal cost depends only on inventory.
These conditions give a price-independent reduced value, not a claim
that the prices have no uncertainty.
\item The finite-penalty problem uses
$G(\mathbf x)=\tfrac12\mathbf x^\top\Phi\mathbf x$, with finite
$\Phi=\Phi^\top\succeq0$. The exact-liquidation problem instead restricts
admissible paths to $\mathbf x(T)=0$. We identify explicitly which
boundary condition is used in each formula.
\item The displayed HJB calculation uses a differentiable value function.
The quadratic solution below supplies that regularity for finite $\Phi$;
for exact liquidation the formulas apply at $t<T$ and the terminal
value is singular away from zero inventory.
\end{enumerate}

\subsection*{Source and reduction}

Use Bergault, Drissi and Gu\'eant \cite{bergault2022}, Eq.~16 with terminal
condition Eq.~17, Remark~2 and Appendix~A. Equation numbers refer to
arXiv:2103.13773v4, 22 March 2022, p.~7.

In the paper's notation, Eq.~16 is
\begin{equation}
\begin{aligned}
0={}&\partial_t\theta
 +\tfrac14(\nabla_q\theta)^\top\eta^{-1}\nabla_q\theta
 +\tfrac12\operatorname{tr}(\Sigma D^2_{SS}\theta)\\
 &-\tfrac\gamma2(q+\nabla_S\theta)^\top\Sigma(q+\nabla_S\theta)
 +(\bar S-S)^\top R^\top(\nabla_S\theta+q).
\end{aligned}
\label{eq:bdg16}
\end{equation}
Set $R=0$. As stated in Remark~2 and Appendix~A of the paper, its quadratic
solution then has no price dependence: $\theta(t,q,S)=q^\top A(t)q$.
Thus $\nabla_S\theta=0$ and $D^2_{SS}\theta=0$, and Eq.~16 becomes
\[
0=\partial_t\theta+\tfrac14(\nabla_q\theta)^\top\eta^{-1}\nabla_q\theta
       -\tfrac\gamma2q^\top\Sigma q.
\]
Translate to this project's cost-minimization convention using
\[
 q=\mathbf x,\qquad \theta=-\mathcal V,\qquad
 \eta=H,\quad\Sigma=C,\quad\gamma=2\lambda,\quad\Gamma=\Phi/2.
\]
Multiplying the reduced equation by $-1$ gives Eq.~\eqref{eq:pde} below.
The paper uses $\dot q=v_{\rm paper}$, whereas we use
$\dot{\mathbf x}=-\mathbf v$. Hence
$v_{\rm paper}^*=\eta^{-1}\nabla_q\theta/2=-\mathbf v^*$, which gives
the feedback in Eq.~\eqref{eq:pde}. Its terminal condition
$\theta(T,q,S)=-q^\top\Gamma q$ becomes
$\mathcal V(\mathbf x,T)=\tfrac12\mathbf x^\top\Phi\mathbf x$.
No discrete-to-continuum derivation is needed for this route.

\subsection*{Classical reference solution}

The corresponding reduced cost-to-go is
\begin{equation}
 \mathcal V_\Phi(\mathbf x,t)=\inf_{\mathbf v(\cdot)}
 \left\{\int_t^T(\mathbf v^\top H\mathbf v+
 \lambda\mathbf x^\top C\mathbf x)\,dr
 +\tfrac12\mathbf x(T)^\top\Phi\mathbf x(T)\right\},
 \qquad\dot{\mathbf x}=-\mathbf v.
 \label{eq:cost}
\end{equation}
For deterministic schedules, Gaussian cost gives a certainty-equivalent risk
term $\gamma\operatorname{Var}/2$, hence $\gamma=2\lambda$.
More explicitly, for $d\mathbf S_r=C^{1/2}d\mathbf B_r$, the random
inventory exposure of a deterministic schedule is
$\int_t^T\mathbf x(r)^\top C^{1/2}d\mathbf B_r$. Its variance is
\[
 \operatorname{Var}\!\left[\int_t^T\mathbf x(r)^\top C^{1/2}d\mathbf B_r\right]
 =\int_t^T\mathbf x(r)^\top C\mathbf x(r)\,dr.
\]
If a Gaussian execution cost $K$ has mean $m$ and variance $\sigma_K^2$,
then $\gamma^{-1}\log\mathbb E[e^{\gamma K}]
=m+\gamma\sigma_K^2/2$ for $\gamma>0$, with its continuous risk-neutral
limit at $\gamma=0$. This explains the risk term for deterministic
schedules. The reduction of the source HJB above provides the
signed-control feedback without assuming that every adapted schedule
produces a Gaussian cost.

\subsection*{Derive the HJB from the Bellman principle}

For a short interval of length $h$, choose a constant rate $\mathbf v$,
then use the optimal continuation. With
$L(\mathbf x,\mathbf v)=\mathbf v^\top H\mathbf v+
\lambda\mathbf x^\top C\mathbf x$, the Bellman principle gives
\[
 \mathcal V(\mathbf x,t)=\inf_{\mathbf v\in\mathbb R^M}
 \left\{hL(\mathbf x,\mathbf v)
 +\mathcal V(\mathbf x-h\mathbf v,t+h)\right\}+o(h).
\]
For each fixed rate, Taylor expansion gives
\[
 \mathcal V(\mathbf x-h\mathbf v,t+h)
 =\mathcal V(\mathbf x,t)+h\mathcal V_t
 -h\mathbf v^\top\nabla\mathcal V+o(h).
\]
Subtract $\mathcal V$, divide by $h$ and let $h\downarrow0$:
\begin{equation}
 0=\mathcal V_t+\min_{\mathbf v\in\mathbb R^M}
 \left\{\mathbf v^\top H\mathbf v
 -\mathbf v^\top\nabla\mathcal V+\lambda\mathbf x^\top C\mathbf x\right\}.
 \label{eq:stage2-bellman}
\end{equation}
The minus sign comes from the selling convention $\dot{\mathbf x}=-\mathbf v$.
The control derivative and Hessian of the Hamiltonian are
\[
 2H\mathbf v-\nabla\mathcal V=0,\qquad 2H\succ0.
\]
Thus its unique global minimizer is
$\mathbf v^*=H^{-1}\nabla\mathcal V/2$. Substitution contributes
$-\tfrac14(\nabla\mathcal V)^\top H^{-1}\nabla\mathcal V$.
The specialized HJB and feedback are
\begin{equation}
 \boxed{\mathcal V_t-\tfrac14(\nabla\mathcal V)^\top H^{-1}\nabla\mathcal V
 +\lambda\mathbf x^\top C\mathbf x=0,\qquad
 \mathbf v^*=\tfrac12H^{-1}\nabla\mathcal V.}
 \label{eq:pde}
\end{equation}
For $g=\nabla\mathcal V$, the identity
\[
 \mathbf v^\top H\mathbf v-\mathbf v^\top g
 =(\mathbf v-\tfrac12H^{-1}g)^\top H(\mathbf v-\tfrac12H^{-1}g)
 -\tfrac14g^\top H^{-1}g
\]
shows that the feedback uniquely minimizes the Hamiltonian.
\subsection*{Derive the Riccati equation}

The costs are quadratic and the dynamics are linear. Try a symmetric
quadratic value
$\mathcal V=\tfrac12\mathbf x^\top P(t)\mathbf x+c(t)$.
This is an ansatz to verify, rather than an additional approximation.
Its derivatives are
\[
 \mathcal V_t=\tfrac12\mathbf x^\top\dot P\mathbf x+\dot c,
 \qquad\nabla\mathcal V=P\mathbf x.
\]
Substitute into the HJB and collect powers of $\mathbf x$:
\[
 0=\tfrac12\mathbf x^\top
 \left(\dot P-\tfrac12PH^{-1}P+2\lambda C\right)\mathbf x+\dot c.
\]
Since this must hold for every inventory, its quadratic matrix and its
constant term vanish separately. Matching the terminal cost gives
\begin{equation}
 \dot P=\tfrac12PH^{-1}P-2\lambda C,\qquad P(T)=\Phi,
 \qquad\dot c=0,\quad c(T)=0,
 \qquad\mathbf v^*=\tfrac12H^{-1}P\mathbf x.
 \label{eq:ric}
\end{equation}
Thus $c=0$ and the optimized initial cost is
$\mathcal V_\Phi(\mathbf X,0)=\tfrac12\mathbf X^\top P(0)\mathbf X$.
Along an admissible path,
\[
 J_\Phi[\mathbf v]-\mathcal V_\Phi(\mathbf X,0)
 =\int_0^T(\mathbf v-\tfrac12H^{-1}P\mathbf x)^\top
 H(\mathbf v-\tfrac12H^{-1}P\mathbf x)\,dt\ge0,
\]
where $J_\Phi$ is the cost in Eq.~\eqref{eq:cost} at $t=0$.
This supplies a direct optimality check.

\subsection*{Derive the exact-liquidation trajectory}

For $\mathbf x(0)=\mathbf X$, $\mathbf x(T)=0$, eliminate the rate by
$\mathbf v=-\dot{\mathbf x}$. The functional to minimize becomes
\[
 J[\mathbf x]=\int_0^T
 (\dot{\mathbf x}^\top H\dot{\mathbf x}
 +\lambda\mathbf x^\top C\mathbf x)\,dt.
\]
Both endpoint variations vanish. The Euler--Lagrange equation is
\[
 \frac{d}{dt}(2H\dot{\mathbf x})-2\lambda C\mathbf x=0,
 \qquad H\ddot{\mathbf x}=\lambda C\mathbf x.
\]
Strict convexity of the derivative term with fixed endpoints makes its
stationary path the unique minimizer. To decouple the equation, diagonalize
$H^{-1/2}CH^{-1/2}=U\operatorname{diag}(\mu_i)U^\top$ and set
$\kappa_i=\sqrt{\lambda\mu_i}$. In coordinates
$\mathbf y=U^\top H^{1/2}\mathbf x$, multiplication of the
Euler--Lagrange equation by $U^\top H^{-1/2}$ gives
\[
 \ddot y_i=\kappa_i^2y_i,\qquad
 y_i(0)=(U^\top H^{1/2}\mathbf X)_i,\qquad y_i(T)=0.
\]
For $\kappa_i>0$ the two elementary solutions are
$e^{\kappa_i t}$ and $e^{-\kappa_i t}$. The terminal condition selects
$y_i(t)=b_i\sinh[\kappa_i(T-t)]$, and the initial condition sets
$b_i=y_i(0)/\sinh(\kappa_iT)$. Transforming back gives
\begin{equation}
 \mathbf x(t)=H^{-1/2}U\operatorname{diag}\!\left(
 \frac{\sinh[\kappa_i(T-t)]}{\sinh(\kappa_iT)}\right)U^\top H^{1/2}\mathbf X.
 \label{eq:traj}
\end{equation}
Differentiating each mode gives
$-\dot y_i=\kappa_i\coth[\kappa_i(T-t)]y_i$. Comparison with
$\mathbf v^*=H^{-1}P\mathbf x/2$ yields the hard-limit feedback matrix
$P(t)=2H^{1/2}U\operatorname{diag}(\kappa_i\coth[\kappa_i(T-t)])U^\top H^{1/2}$.
For zero-risk modes, use
$\sinh(\kappa s)/\sinh(\kappa T)\to s/T$ and
$\kappa\coth(\kappa s)\to1/s$.
Appendix~A of \cite{bergault2022} gives the $R=0$ Riccati solution;
Almgren--Chriss \cite{almgrenchriss2000}, App.~A.3, gives the classical
inventory modes.

\subsection*{Finite terminal penalty and the boundary condition}

For a finite penalty the endpoint is free. Its variation gives
\[
 2H\dot{\mathbf x}(T)+\Phi\mathbf x(T)=0,
 \qquad\mathbf v(T)=\tfrac12H^{-1}\Phi\mathbf x(T).
\]
This condition generally leaves residual inventory; it does not impose
$\mathbf x(T)=0$. A formula valid without covariance/penalty commutation is
\begin{equation}
 \begin{pmatrix}Y(s)\\Z(s)\end{pmatrix}
 =\exp\!\left[s\begin{pmatrix}0&\tfrac12H^{-1}\\2\lambda C&0\end{pmatrix}\right]
 \begin{pmatrix}I\\\Phi\end{pmatrix},\qquad P(T-s)=Z(s)Y(s)^{-1}.
 \label{eq:finite-penalty}
\end{equation}
Indeed, $Y'=H^{-1}Z/2$ and $Z'=2\lambda CY$ give
$(ZY^{-1})'=2\lambda C-(ZY^{-1})H^{-1}(ZY^{-1})/2$ and $P(T)=\Phi$.
This identity is derived here from Eq.~\eqref{eq:ric}. A scalar-mode
finite-penalty solution requires $H^{-1/2}\Phi H^{-1/2}$ and
$H^{-1/2}CH^{-1/2}$ to commute.

\subsection*{Scope and exit criterion}

$R=0$ removes mean reversion, not Brownian price noise. Price independence
uses the paper's terminal data and its $R=0$ solution. A feasible signed-control
optimum also solves the sales-only problem at that initial state; it does
not establish a sales-constrained value function over an entire inventory grid.

Step 2's source reduction is recorded. In Step 3, review its signs, factors
of two and terminal data before accepting the benchmark.

\begin{thebibliography}{9}
\bibitem{almgrenchriss2000} R.~Almgren and N.~Chriss,
\emph{Optimal Execution of Portfolio Transactions}, J.~Risk 3(2), 5--39
(2000), Appendix~A.3.

\bibitem{bergault2022} P.~Bergault, F.~Drissi and O.~Gu\'eant,
\emph{Multi-asset optimal execution and statistical arbitrage strategies
under Ornstein--Uhlenbeck dynamics}, SIAM J.~Financial Math.~13(1),
353--390 (2022). \url{https://arxiv.org/pdf/2103.13773v4}.

\end{thebibliography}
\end{document}
